\documentclass[10pt,a4paper,twocolumn]{article}

% --- Geometry: ACM sigconf-like proportions ---
\usepackage[a4paper,
            top=1.9cm, bottom=2.0cm,
            left=1.7cm, right=1.7cm,
            columnsep=0.6cm]{geometry}

% --- Typography ---
\usepackage[T1]{fontenc}
\usepackage{libertine}        % Linux Libertine for body
\usepackage[libertine]{newtxmath}
\usepackage[scaled=0.85]{beramono}  % SPIR-V code listings
\renewcommand{\sfdefault}{uop}
\usepackage{microtype}

% --- Layout helpers ---
\usepackage{enumitem}
\usepackage{adjustbox}
\setlist{nosep, leftmargin=*, topsep=2pt, partopsep=0pt, parsep=2pt, itemsep=2pt}
\setlength{\parskip}{2pt}
\setlength{\parindent}{10pt}
\widowpenalty=10000
\clubpenalty=10000

% --- Code listings (SPIR-V diff) ---
\usepackage{xcolor}
\definecolor{codebg}{HTML}{F7F7F2}
\definecolor{kwcolor}{HTML}{8B0000}
\definecolor{cmtcolor}{HTML}{6A737D}
\usepackage{listings}
\lstdefinelanguage{spirv}{
  morekeywords={OpLabel,OpULessThan,OpSelectionMerge,OpBranchConditional,
                OpAccessChain,OpLoad,OpBranch,OpPhi,OpTypeFloat,OpTypePointer,
                OpTypeVector,OpTypeFunction,OpVariable,OpFunction,OpFunctionEnd,
                OpReturn,OpStore},
  sensitive=true,
  morecomment=[l]{;},
  morestring=[b]"
}
\lstset{
  language=spirv,
  basicstyle=\ttfamily\scriptsize,
  keywordstyle=\color{kwcolor}\bfseries,
  commentstyle=\color{cmtcolor}\itshape,
  backgroundcolor=\color{codebg},
  frame=single,
  framerule=0.3pt,
  rulecolor=\color{black!25},
  framesep=3pt,
  xleftmargin=4pt,
  xrightmargin=4pt,
  aboveskip=4pt,
  belowskip=2pt,
  breaklines=true,
  showstringspaces=false,
  columns=fullflexible,
  keepspaces=true,
}

% --- Section styling ---
\usepackage{titlesec}
\titleformat{\section}{\bfseries\fontsize{10.5pt}{12pt}\selectfont}{\thesection}{0.5em}{}
\titleformat{\subsection}{\bfseries\fontsize{9.5pt}{11pt}\selectfont}{\thesubsection}{0.5em}{}
\titlespacing*{\section}{0pt}{8pt}{3pt}
\titlespacing*{\subsection}{0pt}{5pt}{2pt}

% --- Hyperref last ---
\usepackage[colorlinks=true,
            linkcolor=blue!60!black,
            citecolor=blue!50!black,
            urlcolor=blue!60!black,
            bookmarks=true,
            bookmarksnumbered=true,
            unicode=true]{hyperref}

% --- Title block (no \maketitle, custom for compact two-column) ---
\newcommand{\papertitle}[1]{%
  \begin{center}
    {\bfseries\fontsize{14pt}{16pt}\selectfont #1\par}
  \end{center}
  \vspace{-2pt}
}
\newcommand{\paperauthor}[1]{%
  \begin{center}
    {\itshape\fontsize{9.5pt}{11pt}\selectfont #1\par}
  \end{center}
  \vspace{2pt}
}

\hypersetup{
  pdftitle={API-Constant-Seeded Bounds-Check Elision for WGSL Compute Shaders in Naga},
  pdfauthor={IICT 2026 Practice Paper Submission},
  pdfsubject={Compiler Optimization, WebGPU, Naga, Interval Analysis},
  pdfkeywords={WebGPU, WGSL, Naga, bounds check elimination, interval analysis, SPIR-V}
}

\begin{document}

\twocolumn[
\papertitle{API-Constant-Seeded Bounds-Check Elision for\\WGSL Compute Shaders in Naga}
\paperauthor{Practice Paper Submission --- IICT 2026: Innovations in Compiler Technology\\
Indian Institute of Science, Bengaluru \quad\textbar\quad October 2--3, 2026}
\begin{adjustbox}{minipage=0.92\textwidth, center}
\fontsize{9pt}{11pt}\selectfont
\noindent\textbf{Abstract.}\enspace
The WebGPU specification mandates that out-of-bounds array, vector, and matrix
accesses be redirected to a safe location rather than left as undefined
behaviour. Naga, the shader translator at the core of the \texttt{wgpu}
ecosystem, enforces this through the \texttt{BoundsCheckPolicy} framework in
\texttt{naga/src/proc/index.rs}. However, Naga's elision logic---the function
\texttt{access\_needs\_check}---only folds indices that resolve to compile-time
constants via \texttt{try\_resolve\_to\_constant}, ignoring the value ranges the
WebGPU specification guarantees for built-in inputs such as
\texttt{local\_invocation\_index}. As a consequence, the canonical compute idiom
\texttt{tile[local\_invocation\_index]}, where \texttt{tile} is sized exactly to
\texttt{@workgroup\_size}, unconditionally pays a runtime bounds check that
fragments the SPIR-V into a divergent branch-and-merge block. This work
integrates a lightweight interval analysis pass into Naga's validator, seeding
\texttt{u32} intervals from \texttt{@workgroup\_size} and propagating them
through constant addition and subtraction using strictly checked unsigned
arithmetic. The pass eliminates redundant control flow across all backends
without requiring backend modifications, targets the open \texttt{gfx-rs/wgpu}
issue~\#6345, and is positioned relative to the GSOHC synchronization-optimization
lineage and to LLVM's Inductive Range Check Elimination.
\end{adjustbox}
\vspace{8pt}
]

\section{Motivation}

Naga relies on \texttt{access\_needs\_check} in \texttt{naga/src/proc/index.rs}
to decide whether an array, vector, or matrix access requires a runtime guard.
Under the \texttt{ReadZeroSkipWrite} policy---the default for native
\texttt{wgpu} deployments---failing to elide a check produces SPIR-V containing
\texttt{OpULessThan}, \texttt{OpSelectionMerge}, \texttt{OpBranchConditional},
and a merge block terminating in \texttt{OpPhi}. The result is a four-instruction
divergent control-flow diamond surrounding every dynamic access, plus the
fragmentation of an otherwise straight-line basic block into three.

Although the WebGPU specification~\cite{webgpu} guarantees that
\texttt{local\_invocation\_index} is bounded by the pipeline-time constant
\texttt{@workgroup\_size}, Naga treats it as an opaque runtime value. Indexing
\texttt{var<workgroup>} memory by \texttt{local\_invocation\_index} therefore
injects divergent control flow into latency-bound compute kernels---reductions,
prefix scans, tiled convolutions---where branch divergence is highly
detrimental to warp occupancy and instruction-level parallelism. The goal of
this work is to close this gap using information already available during
shader validation, without extending the WebGPU API surface or modifying any
backend code generator.

\section{Interval Propagation and Soundness}

Because \texttt{access\_needs\_check} already receives a
\texttt{\&valid::FunctionInfo} parameter, the hook point for a smarter answer
exists without changing function signatures. We extended \texttt{ExpressionInfo}
in \texttt{naga/src/valid/analyzer.rs} with a new field
\texttt{range: Option<(u32, u32)>}, computed during the existing single forward
walk over the expression arena. This design follows the classical abstract
interpretation framework of Cousot \& Cousot~\cite{cousot1977}, where the
lattice is the flat product domain over $\texttt{Option}((\mathbb{U}_{32})^2)$:
the bottom element is \texttt{None} (unknown range, fall back to runtime check)
and the top elements are concrete $[\text{lo}, \text{hi}]$ intervals.

The primary challenge is soundness. WGSL~\cite{wgsl} \texttt{u32} arithmetic
wraps modulo $2^{32}$ rather than producing negative numbers; using signed
interval arithmetic would silently describe different values than the hardware
computes. For example, \texttt{0u - 1u} evaluates to \texttt{4294967295u}, not
$-1$. A naive signed interval $[0, 63] - 1 = [-1, 62]$ would falsely claim that
thread~0 computes $-1$ on the index path, when in fact it produces the maximum
\texttt{u32} value and triggers a runtime clamp at best, or---if the analysis
had been used to \emph{elide} the clamp---an out-of-bounds memory access at
worst.

To prevent false positives, the propagation strictly uses \texttt{u32} bounds
with \texttt{u32::checked\_add} and \texttt{u32::checked\_sub}. If an operation
would underflow or overflow (for instance, \texttt{lid - 1u} where
\texttt{lid}~$\in [0, 63]$), the interval safely collapses to \texttt{None},
and Naga falls back to its existing conservative runtime check. This ensures
the only failure mode is a missed optimisation, never an unsound elision---the
classical precision-versus-soundness trade-off of abstract interpretation,
resolved here unambiguously in favour of soundness.

\section{Bounds Check Resolution}

During the \texttt{Expression::FunctionArgument} arm of
\texttt{process\_expression}, the pass consults the argument's
\texttt{Binding::BuiltIn} attribute. If it is a \texttt{LocalInvocationIndex}
or \texttt{LocalInvocationId}, the range is seeded from the enclosing entry
point's statically known \texttt{workgroup\_size}. For binary expressions,
intervals are propagated only when the opposing operand is a compile-time
constant; this restriction keeps the analysis within $O(|E|)$ where $|E|$ is the
expression arena size, and avoids the need for a fixpoint iteration.

In \texttt{access\_needs\_check}, a single disjunct was added: if the index is
a \texttt{GuardedIndex::Expression} and \texttt{info[expr].range} evaluates to
\texttt{Some((lo, hi))} with \texttt{hi < length}, the function returns
\texttt{None}, indicating no check is needed. Because this logic lives in the
shared \texttt{proc/index.rs} module---already consulted by every backend---the
optimisation is automatically inherited by the SPIR-V, MSL, and HLSL backends
with zero changes to backend call sites or signatures. The patch comprises
approximately 50 lines of core logic.

\section{Architectural Overview}

Figure~\ref{fig:diff} shows the SPIR-V emitted by Naga for the canonical idiom
\texttt{tile[lid]} where \texttt{lid} is \texttt{local\_invocation\_index} and
\texttt{tile} is \texttt{array<f32, 64>} with \texttt{@workgroup\_size(64)}.
Instruction counts per line are annotated as \texttt{; 1}.

\begin{figure}[t]
\begin{lstlisting}
; --- Before (ReadZeroSkipWrite policy) ---
%header     = OpLabel                                  ; 1
%cmp        = OpULessThan %bool %lid %64               ; 2
            OpSelectionMerge %merge None               ; 3
            OpBranchConditional %cmp %in_bounds %merge ; 4
%in_bounds  = OpLabel
%ptr        = OpAccessChain %ptr_ty %tile %lid
%val        = OpLoad %f32 %ptr
            OpBranch %merge                           ; 5
%merge      = OpLabel
%result     = OpPhi %f32 %zero %header %val %in_bounds ; 6
; 6 wasted instructions, 3 basic blocks, divergent

; --- After (interval proves lid in [0, 63]) ---
%ptr        = OpAccessChain %ptr_ty %tile %lid
%val        = OpLoad %f32 %ptr
; 0 wasted instructions, 1 basic block, straight-line
\end{lstlisting}
\caption{SPIR-V generated by Naga for \texttt{tile[lid]} before and after
API-constant-seeded elision. Six instructions and two extra basic blocks are
removed; the access becomes straight-line and eligible for backend
instruction scheduling.}
\label{fig:diff}
\end{figure}

The eliminated instructions survive into AMDGCN and PTX because WebGPU's
robustness semantics forbid the downstream driver from removing them---the
runtime clamp is a semantic guarantee, not an optimiser hint. Eliminating theclamp at the IR level is therefore the only sound point at which the cost can be avoided.

\section{Supplementary Optimization: Barrier Elision for workgroupUniformLoad}
\label{sec:barrier-elision}

The WebGPU specification~\cite{webgpu} requires that
\texttt{workgroupUniformLoad} be sequenced with respect to prior and
subsequent workgroup-level memory operations.  Naga's SPIR-V backend enforces
this by emitting both a pre-barrier and a post-barrier
(\texttt{OpMemoryBarrier} with \texttt{WORKGROUP} semantics) around each
\texttt{WorkGroupUniformLoad} statement.  For shaders that issue consecutive
\texttt{workgroupUniformLoad} calls, or patterns of the form
\texttt{workgroupBarrier(); let x = workgroupUniformLoad(\&ptr);}, the
resulting barriers are redundant and contribute measurable instruction overhead.

We introduce two lightweight forward-scan optimizations in the SPIR-V backend
that identify and elide these redundant barriers at compile time, without
altering Naga's IR, validator, or other backends.

\subsection{Track~3: Post-Barrier Elision}

When the SPIR-V backend processes a \texttt{WorkGroupUniformLoad} statement,
it scans the remaining statements in the current block for the next
synchronization point---either an explicit \texttt{Barrier(WORK\_GROUP)} or
another \texttt{WorkGroupUniformLoad} whose pre-barrier can serve as the
post-barrier for the preceding load.  If such a point is found in the same
block, the post-barrier is omitted.

This targets the common reduction idiom where two or more
\texttt{workgroupUniformLoad} calls appear back-to-back inside a loop, or a
\texttt{workgroupBarrier()} immediately follows the WGUL.  A representative
kernel compiled under the \texttt{ReadZeroSkipWrite} policy shows a static
saving of 8 SPIR-V words (2 WORK\_GROUP barriers, each 4 words) for
the \texttt{wgul\_reduce} test shader: 699 words fall to 691.

\subsection{Track~3b: Pre-Barrier Elision}

Symmetrically, when the backend encounters an explicit
\texttt{Barrier(WORK\_GROUP)} statement, it scans forward for the first
non-\texttt{Emit} statement.  If that statement is a
\texttt{WorkGroupUniformLoad}, the barrier is elided---the WGUL's pre-barrier
already provides the required visibility guarantee.

This targets patterns where an explicit \texttt{workgroupBarrier()} directly
precedes a \texttt{workgroupUniformLoad}.  Such sequences occur when shaders
alternate between uniform (WGUL) and non-uniform workgroup memory accesses.
The \texttt{track3b\_test} shader---which contains two
\texttt{Barrier}$\rightarrow$\texttt{WGUL} pairs and one consecutive
\texttt{WGUL}$\rightarrow$\texttt{WGUL} sequence---saves 8 SPIR-V words
(582 $\rightarrow$ 574) when Track~3b is active.

\subsection{Benchmark Results}

We measured the timing impact on an NVIDIA GeForce RTX~4060 Laptop GPU using a
dedicated \texttt{wgpu} benchmark harness.  The harness dispatches the
\texttt{track3b\_test} shader with GPU timestamp queries for sub-microsecond
resolution, 10 timed passes per configuration, and bootstrap 95\% confidence
intervals on the median.  Dispatches per pass are scaled by problem size
(10,000 for 64K through 10 for 16M) to maintain consistent pass durations.

\begin{table}[t]
\centering
\caption{Track~3b timing impact for the \texttt{track3b\_test} kernel under
the \texttt{ReadZeroSkipWrite} bounds-check policy.  Word counts are per
dispatch; durations are median GPU compute-pass elapsed times.}
\label{tab:track3b-timing}
\small
\begin{tabular}{lcccr}
\hline
\textbf{Size} & \textbf{Without T3b (\textmu s)} & \textbf{With T3b (\textmu s)} & \textbf{Delta} & \textbf{Words} \\
\hline
64K   & 2.594 & 2.520 & $-$2.9\% & 582 $\rightarrow$ 574 \\
1M    & 22.420 & 21.395 & $-$4.6\% & 582 $\rightarrow$ 574 \\
10M   & 258.561 & 242.746 & $-$6.1\% & 582 $\rightarrow$ 574 \\
16M   & 385.606 & 363.379 & $-$5.8\% & 582 $\rightarrow$ 574 \\
\hline
\end{tabular}
\end{table}

Table~\ref{tab:track3b-timing} summarises the results.  Track~3b yields a
consistent 2.9--6.1\% wall-clock speedup across all four problem sizes.  The
control kernel (\texttt{wgul\_reduce}, which does not contain the
\texttt{Barrier}$\rightarrow$\texttt{WGUL} pattern) shows no significant
timing change (0.3--3.4\% variance, consistent with GPU run-to-run noise).

For Track~3 (post-barrier elision in isolation), the static SPIR-V word-count
saving is confirmed (8 words, 2 barriers) but single-run GPU timing noise
prevents a statistically significant wall-clock measurement at the
benchmark's current sample count.  The direction is consistently positive
across problem sizes and policies, consistent with the SPIR-V-level reduction.

Both optimisations are implemented in \texttt{naga/src/back/spv/block.rs} as
targeted forward scans in the statement-emitting loop, adding approximately
35 lines of Rust in total.  Neither requires changes to Naga's IR, validator,
or other backends.

\section{Current Status and Future Work}

At the time of writing, the bounds-check elision (Track~2) and the barrier-elision
optimisations (Track~3 and Track~3b, \S\ref{sec:barrier-elision}) are complete
on a feature branch and undergoing testing against the Naga test suite; a pull
request targeting \texttt{gfx-rs/wgpu} issue~\#6345~\cite{issue6345} is in
preparation for upstream review.  Preliminary analysis confirms the bounds-check
elision eliminates the branch-and-merge overhead for the
\texttt{ReadZeroSkipWrite} path, yielding expected runtime improvements in
low-occupancy, latency-bound kernels.  As anticipated, the \texttt{Restrict}
policy---which emits only a single \texttt{UMin} instruction rather than a full
diamond---shows negligible wall-clock improvement, because GPU warp schedulers
effectively hide this single-cycle ALU dependency.

Remaining work involves finalising negative test cases for wraparound scenarios
(\texttt{lid - 1u} on thread 0, multiplicative overflow) and benchmarking on
three vendor ISAs (AMD RDNA3, NVIDIA Ada, Intel Arc) to confirm that the static
SPIR-V instruction-count reduction translates into runtime gains across driver
stacks.  The barrier-elision optimisations (Track~3/3b) will also be tested on
multi-vendor hardware to confirm they are equally effective across GPU
architectures.  Future extensions include propagating intervals through
multiplicative index expressions such as \texttt{lid * stride} (requires
divisibility reasoning), integrating \texttt{global\_invocation\_id}
ranges---the latter requires knowing \texttt{numWorkgroups} at dispatch time, a
value not available at shader compilation and therefore demanding a
deferred-validation or specialisation-constant mechanism---and extending the
barrier-elision forward scan to track state across nested control flow.  We also
plan to compare the implementation against LLVM's IRCE pass~\cite{llvm-irce} on
equivalent SYCL kernels, to quantify the gap between Naga's bespoke IR and the
LLVM middle-end pipeline that vendor-agnostic HPC runtimes (SYCL, Julia
\texttt{CUDA.jl}, OpenMP target) inherit for free.

\section*{Acknowledgements}

The author thanks the \texttt{gfx-rs/wgpu} maintainers for the accessible
codebase and the IICT 2026 program committee for the practice-paper track that
motivated this write-up.

\footnotesize
\section*{References}
\begin{enumerate}[leftmargin=*, label={\arabic*.}]
\item gfx-rs/wgpu repository. \url{https://github.com/gfx-rs/wgpu}
\item \texttt{gfx-rs/wgpu} Issue \#6345, ``[naga spv-out] Bounds check generation could be cleaner and smarter.''
\url{https://github.com/gfx-rs/wgpu/issues/6345}\label{issue6345}
\item W3C. \emph{WebGPU Specification}. \url{https://www.w3.org/TR/webgpu/}\label{webgpu}
\item W3C. \emph{WGSL Specification}. \url{https://www.w3.org/TR/WGSL/}\label{wgsl}
\item Cousot, P. \& Cousot, R. ``Abstract interpretation: a unified lattice model for static analysis of programs by construction or approximation of fixpoints.'' \emph{POPL}, 1977.\label{cousot1977}
\item Khronos Group. \emph{Vulkan Specification (SPIR-V environment)}. \url{https://www.khronos.org/registry/vulkan/specs/latest/html/vkspec.html}
\item LLVM Project. ``Inductive Range Check Elimination (IRCE) pass.'' \url{https://llvm.org/devmtg/2015-04/slides/Hal_Finkel_IRCE.pdf}\label{llvm-irce}
\end{enumerate}

\end{document}
